\subsection{Positive diagonal pairings and analytic matrix orders}
\label{sec:v19-comparison}
The finite square necessity-and-sufficiency statement has an exact
predecessor, not merely a common determinant identity. For
$X=\{x_1,\ldots,x_m\}$, let
$A_{ik}=v_i(x_k)$ and $B_{jk}=\phi_j(x_k)$, with both matrices of
row rank $d$. The pairing is $A\operatorname{diag}(w)B^t$.
Banaji--Pantea \cite[Definition 2.24 and Lemma 2.36]{BanajiPantea2016}
characterize the corresponding common weak sign, with a nonzero
product of maximal minors, by rank preservation under positive
diagonal scalings. Their criterion applies because
\[
 \rank(AD_1B^tD_2A)=d
       \quad\Longleftrightarrow\quad\det(AD_1B^t)\ne0;
\]
$D_2A$ is onto. Positive weights and positive probability weights
are equivalent by scalar normalization. The compact-space version
uses full-support measures; posterior normalization and unconditional
history mass are the additional experiment-specific conclusions.
Andr\'eief's identity alone would not attribute this necessity result.

For rectangular evaluations $A\in\R^{d\times m}$ and
$B\in\R^{k\times m}$, full column rank for every positive $w$ is
equivalent to
\[
 \ker A\cap\operatorname{diag}(w)\operatorname{im}B^t=\{0\}
                     \quad\hbox{for every }w>0.
\]
The finite positive-diagonal/sign-vector alternative is the one in
M\"uller et al.\ \cite[Lemma 2.1]{MullerEtAl2016}. Our product-cone
formulation does not claim a new finite sign-vector theorem. Its
compact-space rank alternatives specify full-support witnesses for
entire kernels, and its quantitative form returns to the actual
history map. Proposition~\ref{prop:v19-pentagon} shows why preserving
only a fixed square subspace containing evidence would lose part of
this experiment class.

Similarly, the Smith form and its singular-value orders are classical
matrix facts; see Kaveh--Makhnatch \cite[Theorem 1.1]{KavehMakhnatch}.
They are not themselves a memory theorem. Here the affine acquisition
law supplies constants uniform through rank loss, while
Theorem~\ref{thm:v19-covariance-realization} realizes arbitrary small
matrix germs inside one fixed positive physical experiment. Together
these give all joint memory exponents in
Corollary~\ref{cor:v19-memory-phases}, including priors uniformly
equivalent to a fixed full-support measure. A formal matrix
factorization alone supplies neither physical realization nor a risk
comparison uniform in the entire integer budget.
